\documentclass[planificacion.tex]{subfiles}
\begin{document}
  \begin{longtable}{|p{3cm}|p{13cm}|p{5cm}|p{5cm}|}  \hline
         {\centering EJE TEMÁTICO} &  {\centering SABERES} &  {\centering ESTRATEGIAS DE ENSEÑANZA-APRENDIZAJE} & {\centering EVALUACIÓN} \\ \hline   \endhead
FORMAS DE ENERGÍA Y SUS TRANSFORMACIONES. &
    \begin{itemize} 
        \item Conocer los principios que
rigen las transformaciones
energéticas. 
        \begin{itemize}
            \item Interpretación de las distintas escalas de temperatura.
\item Reconocimiento de la fuerza impulsora (diferencia de
temperatura) que permite la transferencia de calor.
\item Diferenciación de los mecanismos de transferencia de
calor.(conducción, convección, radiación)
\item Análisis comparativo de dilataciones lineales, superficiales
y cúbicas. 
\item Identificación del Principio Cero de la Termodinámica, y su
aplicación en Calorimetría.
\item Elaboración de la relación existente entre balance térmico,
capacidad calorífica y calor específico.
\item Descripción de los fenómenos de dilatación y deformación
en materiales de uso industrial.
        \end{itemize}
    \end{itemize}
    & \estrategias & \evaluacion \\ \hline
   PRINCIPIOS FUNDAMENTALES DE LA TERMODINÁMICA. & 
    \begin{itemize}
        \item Analizar el trabajo en un
sistema cerrado adiabático.
\begin{itemize}
    \item Conceptualización de las características de trabajo
termodinámico.
\item Identificación de las condiciones de un sistema cerrado.
\item Caracterización del principio de la accesibilidad
adiabática.
\item Caracterización de los principios enunciados por Lord
Kelvin y Clausius.
\item Identificación del principio de conservación de la energía
entre estados de equilibrio.
\item Conocimiento de las interacciones másica, mecánica y
térmica de modo sistémico.
\item Representación gráfica de transformaciones isotérmicas,
isobáricas, isocóricas, adiabáticas, politrópicas.
\item Conceptualización de la energía interna como variable de
estado.
\item Análisis crítico del Primer Principio de la Termodinámica.
\item Caracterización de trabajo de expansión y circulación.
\item Representación gráfica de los distintos tipos de trabajos
termodinámicos.
\end{itemize}
\end{itemize}
 & \estrategias & \evaluacion \\ \hline
 & \begin{itemize}
        \item Conocer las causas y consecuencias del incremento de la Entropía en el Universo.
        \begin{itemize}
            \item  Caracterización de la Entropía de modo universal e industrial.
\item Explicitación de la irreversibilidad en procesos
termodinámicos.
\item Reconocimiento de teorías de cantidad de energía no
utilizable en los sistemas.
\item Análisis de los procesos de transformación de materia y
energía desde una perspectiva de creación-destrucción.
\item Análisis crítico del Segundo Principio de la termodinámica
desde la visión de Clausius y de lo desarrollado por
Kelvin-Planck.
\item Reconocimiento de la máquina ideal de Carnot
contemplando el rendimiento idealista y su reversibilidad.
\item Análisis de diagramas gráficos relacionados con procesos
entrópicos.
\item Formulación del Segundo Principio de la Termodinámica desde una perspectiva técnica - industrial. 
        \end{itemize} 
    \end{itemize}
    
    & \estrategias & \evaluacion \\ \hline
           MÁQUINAS DE COMBUSTIÓN INTERNA & 
    \begin{itemize}
        \item Analizar críticamente la matriz tecnológica aplicada en las máquinas de combustión interna en el mundo energético actual.
    \begin{itemize}
        \item Conocimiento de los aportes sobre motores de
combustión interna de NikolausAugust Otto.
\item Reconocimiento de los tiempos o ciclos termodinámicos
en los motores de combustión interna.
\item Interpretación gráfica de las variables intervinientes en un
ciclo ideal de combustión interna.
\item Análisis comparativo entre los desarrollos tecnológicos de
Nikolaus Otto y Rudolf Karl Diesel, considerando
beneficios y cualidades de los ciclos de motores
generados por ellos.
\item Diferenciación de rendimientos de los distintos ciclos de
combustión interna.
\item Conceptualización de las fórmulas técnicas vinculadas
con las máquinas de combustión interna.
    \end{itemize}
    \end{itemize}
        & \estrategias & \evaluacion \\ \hline
        CICLOS TERMODINÁMICOS DE VAPOR Y GAS. &
    \begin{itemize}
        \item Reconocer y analizar las transformaciones termodinámicas en los ciclos de vapor.
        \begin{itemize}
            \item Caracterización de las transformaciones de un sistema
gaseoso.
\item Interpretación del fenómeno de vaporización, a través de
diagramas, representación, con usos de tablas de vapor.
\item Conceptualización de la Entalpia del líquido y del vapor
para el empleo tecnológico en máquinas refrigerantes.
\item Comprensión de los fenómenos de vaporización e
interpretación de las tablas de vapor.
\item Conocimiento de los principios de Carnot relacionados
con los ciclos de vapor.
\item Reconocimiento de las características de diagramas,
título de vapor, empleando los aportes técnicos de Mollier.
\item Caracterización de los principios desarrollados por William
John Rankine en los ciclos de vapor
\item Reconocimiento de los efectos de presión y temperatura
en los ciclos de vapor. 
        \end{itemize}
        \item Conocer los principios vinculados con los ciclos de gas. 
        \begin{itemize}
            \item Caracterización del funcionamiento de una turbomáquina
motora, con empleo de gas como fluido actuante.
\item Análisis comparativo de las propiedades técnico-térmicas
de los ciclos de vapor y gas.
\item Conceptualización de ciclos de gas de potencia y de
refrigeración según Brayton.
\item Exploración de la fluidez y la presión del aire en un ciclo
termodinámico.
\item Interpretación de las funciones tecnológicas específicas,
del compresor, turbina y otros dispositivos acoplados en la
generación eléctrica.

        \end{itemize}
    \end{itemize}
         & \estrategias & \evaluacion\\ \hline
    \end{longtable}
\end{document}